% ─── Capítulo: Metodología — Integración con Logisim ──────────────────────────
\chapter{METODOLOGÍA — Integración con Logisim Evolution}

\section{Visión General del Proceso}

El compilador TPP no se detiene en la generación del archivo \texttt{output.s} (ensamblador RV32I en formato texto). Para que el programa pueda \textbf{ejecutarse realmente en el procesador} implementado en Logisim Evolution, es necesario un paso adicional: convertir el código ensamblador a \textbf{código máquina hexadecimal} cargable directamente en la memoria RAM del circuito.

Este proceso es realizado por el script \texttt{compile\_to\_logisim.py}, derivado y adaptado del repositorio público \textbf{riscv-asm} (disponible en GitHub), un ensamblador minimalista de RISC-V en Python diseñado específicamente para la carga de programas en simuladores de hardware educativos.

El flujo completo, desde el código fuente hasta la ejecución en el procesador, es el siguiente:

\begin{figure}[H]
\centering
\begin{tikzpicture}[node distance=0.6cm, auto,
    step/.style={rectangle, draw=tppblue, fill=tppblue!8,
                 text width=5.5cm, text centered, rounded corners,
                 minimum height=0.9cm, font=\small\bfseries},
    tool/.style={font=\footnotesize\itshape, text=gray},
    arrow/.style={thick,->,>=stealth,color=tppblue}]

    \node [step] (src)   {Código Fuente (.to)};
    \node [step] (comp)  [below=of src]  {Compilador TPP (C)};
    \node [step] (asm)   [below=of comp] {output.s (RV32I ASM texto)};
    \node [step] (py)    [below=of asm]  {compile\_to\_logisim.py};
    \node [step] (hex)   [below=of py]   {output.hex (Hexadecimal)};
    \node [step] (log)   [below=of hex]  {Procesador Logisim (RAM)};

    \node [tool, right=1.5cm of comp] {\texttt{./tpp archivo.to}};
    \node [tool, right=1.5cm of py]   {\texttt{python3 compile\_to\_logisim.py}};
    \node [tool, right=1.5cm of log]  {Carga manual en RAM};

    \draw [arrow] (src) -- (comp);
    \draw [arrow] (comp) -- (asm);
    \draw [arrow] (asm) -- (py);
    \draw [arrow] (py) -- (hex);
    \draw [arrow] (hex) -- (log);
\end{tikzpicture}
\caption{Pipeline completo hasta ejecución en Logisim}
\label{fig:pipeline-logisim}
\end{figure}

\section{El Script \texttt{compile\_to\_logisim.py}}

\subsection{Origen y Adaptación}

El script fue extraído y adaptado del repositorio \textbf{\texttt{riscv-asm}} disponible en GitHub. Este repositorio implementa un ensamblador RISC-V minimalista en Python, capaz de traducir código ensamblador RV32I a su representación binaria y hexadecimal, orientado a contextos educativos donde no se desea depender de toolchains complejos como \texttt{riscv64-unknown-elf-gcc}.

La adaptación realizada para TPP incluye:
\begin{itemize}
    \item Soporte para el formato de memoria de Logisim Evolution (\texttt{v2.0 raw}).
    \item Manejo de etiquetas relativas (pseudo-instrucción \texttt{j}, \texttt{call}).
    \item Soporte de la sección \texttt{.data} para cadenas de texto literales.
    \item Traducción de la pseudo-instrucción \texttt{li} a su secuencia \texttt{lui + addi}.
\end{itemize}

\subsection{Funcionamiento Interno}

El script realiza los siguientes pasos en orden:

\begin{enumerate}
    \item \textbf{Primera pasada — Recolección de etiquetas:}
    Lee el archivo \texttt{output.s} completo. Para cada línea que termina en \texttt{:} (ej. \texttt{main:}, \texttt{.L0:}), calcula la dirección de memoria correspondiente (contando words de 32 bits desde la dirección base \texttt{0x0000}) y registra la etiqueta en un diccionario de símbolos.

    \item \textbf{Segunda pasada — Ensamblado:}
    Para cada instrucción reconocida, aplica la codificación de la ISA RV32I (formato R, I, S, B, U, J) para producir el \textbf{word de 32 bits} en hexadecimal.
    \begin{itemize}
        \item Instrucciones de tipo R (\texttt{add}, \texttt{sub}, \texttt{and}): campos \texttt{funct7|rs2|rs1|funct3|rd|opcode}.
        \item Instrucciones de tipo I (\texttt{lw}, \texttt{addi}): campo inmediato de 12 bits en complemento a dos.
        \item Instrucciones de tipo S (\texttt{sw}): el inmediato se divide entre \texttt{inst[11:5]} e \texttt{inst[4:0]}.
        \item Instrucciones de tipo B (\texttt{beqz}, \texttt{bne}): el inmediato codifica un offset en palabras con codificación especial.
        \item Pseudo-instrucciones (\texttt{li}, \texttt{mv}, \texttt{j}, \texttt{call}, \texttt{ret}): se expanden a una o más instrucciones reales.
    \end{itemize}

    \item \textbf{Generación del archivo \texttt{output.hex}:}
    El archivo resultante inicia con la cabecera \texttt{v2.0 raw} (requerida por Logisim), seguido de un word hexadecimal por línea:

\begin{verbatim}
v2.0 raw
00000097
00000113
ff410113
...
\end{verbatim}

    \item \textbf{Carga en Logisim:}
    El archivo \texttt{output.hex} se carga en la componente \texttt{RAM} del circuito de Logisim mediante \textit{Right Click $\rightarrow$ Load Image...}. El procesador entonces lee las instrucciones desde la dirección \texttt{0x0000} al iniciar.
\end{enumerate}

\subsection{Codificación de Ejemplo}

Para ilustrar el proceso, la instrucción \texttt{addi sp, sp, -32} se codifica de la siguiente manera en formato I:

\begin{table}[H]
\centering
\begin{tabular}{ccccc|c}
\toprule
\textbf{imm[11:0]} & \textbf{rs1} & \textbf{funct3} & \textbf{rd} & \textbf{opcode} & \textbf{Hex} \\
\midrule
\texttt{111111100000} & \texttt{00010} & \texttt{000} & \texttt{00010} & \texttt{0010011} & \texttt{fe010113} \\
(= -32 en 12 bits) & (sp=x2) & (ADDI) & (sp=x2) & (OP-IMM) & \\
\bottomrule
\end{tabular}
\caption{Codificación de \texttt{addi sp, sp, -32} en RV32I formato I}
\label{tab:codificacion-I}
\end{table}

\section{Arquitectura RISC-V y el Procesador Logisim}

\subsection{El Procesador de Ciclo Único}

El procesador implementado en Logisim Evolution es un \textbf{procesador de ciclo único} (Single-Cycle Processor) para la ISA RISC-V RV32I. Sus componentes principales son:

\begin{itemize}
    \item \textbf{Program Counter (PC):} Registro de 32 bits que apunta a la instrucción actual en RAM.
    \item \textbf{Memoria de Instrucciones (RAM):} Donde se carga el archivo \texttt{output.hex}.
    \item \textbf{Banco de Registros:} 32 registros de 32 bits (x0--x31). x0 siempre es cero.
    \item \textbf{ALU:} Unidad Aritmético Lógica que ejecuta operaciones según el campo \texttt{funct3/funct7}.
    \item \textbf{Unidad de Control:} Decodifica el opcode y genera las señales de control para el datapath.
    \item \textbf{Memoria de Datos (RAM):} Accedida con instrucciones \texttt{lw}/\texttt{sw}.
    \item \textbf{TTY Output:} Dirección mapeada a memoria (\texttt{0xFF000004}) que actúa como terminal de salida de caracteres ASCII.
\end{itemize}

\subsection{Ciclo de Ejecución por Instrucción}

Cada instrucción se ejecuta en un único ciclo de reloj mediante los siguientes pasos del datapath:

\begin{enumerate}
    \item \textbf{Fetch (IF):} Leer la instrucción de 32 bits desde RAM[\texttt{PC}].
    \item \textbf{Decode (ID):} La unidad de control decodifica el opcode. Se leen los registros fuente.
    \item \textbf{Execute (EX):} La ALU realiza la operación aritmética o calcula la dirección de memoria.
    \item \textbf{Memory (MEM):} Si es \texttt{lw}/\texttt{sw}, se accede a la RAM de datos.
    \item \textbf{Write Back (WB):} El resultado se escribe en el registro destino.
    \item \textbf{PC Update:} PC $\leftarrow$ PC + 4 (o PC + offset para saltos).
\end{enumerate}
